\section{Discussion}
\label{sec:discussion}

\paragraph{Position against bespoke per-protocol evaluation.}
The value of the benchmark is comparative, and it is worth stating exactly what it
adds over the status quo. Each timing monitor and protocol we cite carries a real
evaluation of its own: Gioia's timing protection level \citep{gioia2023} is scored
against GNSS scenarios, Network Time Security \citep{rfc8915} and Roughtime
\citep{roughtime} against cryptographic threat models, IEEE~802.1ASed
\citep{ieee8021ased} against a lying-source selection model, and the composed
protection level \citep{baweja2026cti} against a holdover metrology series. These
evaluations are not wrong; they are \emph{incommensurable}. A benchmark's
contribution is to make them commensurable: the same fault menu, the same scoring
rule, the same treatment of the undetectable, applied to whichever monitor one can
express as a protection-level function $\PL(\cdot)$. We do not surpass any of these
results on its own ground---we cannot, and do not claim to---we provide the common
ground on which they can be laid side by side.

The closest neighbours to what we build are recent timing-security evaluation
frameworks, and it is worth saying precisely how this benchmark differs from them.
BAND \citep{ren2025band}, whose name announces a benchmark, is a physics-based
analysis of non-reciprocal (asymmetric) delay-attack \emph{detection} on a single
$100$-km dual-comb optical two-way link; TIMESAFE \citep{groen2024timesafe} is a
machine-learning \emph{detector} for spoofing and replay on the Precision Time
Protocol in a 5G fronthaul; and PTPsec \citep{finkenzeller2024ptpsec} is a PTP
security scheme that detects delay attacks by cyclic path-asymmetry analysis. Each
is a real and valuable result, but each evaluates \emph{one detector on one
protocol}: it asks how well a particular mechanism catches a particular attack on a
particular link. None of them scores an \emph{arbitrary} monitor's \emph{stated}
integrity risk against a \emph{common, fixed} fault menu through the Stanford
integrity diagram, and none encodes the symmetric-delay impossibility as a
structural invariant that forbids crediting detection of the undetectable. The
distinction is not one of quality but of kind: they are monitors, or detector
evaluations, on the object side of the ruler; this paper is the ruler.

\paragraph{What the benchmark deliberately does not claim.}
The honesty-immunity of Section~\ref{sec:immunity} cuts both ways: the benchmark
cannot overclaim, but it also cannot make several claims one might wish it could. It
does not assert that any monitor is safe, that any protection level is tight, that
any fault rate is realistic, or that passing the benchmark implies fitness for any
application. It emits tallies against a fixed rule; interpreting those tallies for a
particular application---choosing the alert limit, judging whether the menu covers
the deployment's real faults, deciding whether the stated integrity risk is
credible---remains the user's responsibility. In particular, a monitor that passes
the entire menu has demonstrated coverage \emph{of that menu}, not of the world.

\paragraph{The honest novelty bound.}
Stated plainly, the novelty is three things and no more. First, the existence of
the open, versioned, cross-protocol conformance benchmark itself---the common
ruler---which, to our knowledge, no predecessor provides; we phrase this ``to our
knowledge'' and take care never to claim a first secure-time protocol, a first
time-RAIM, or the non-existence of timing protection levels, all of which exist and
all of which this benchmark scores rather than replaces. Second, the
verdict-algebra impossibility encoding: making Mizrahi's symmetric-delay
undetectability a type-level invariant, so that reporting an undetectable fault as
detected is unrepresentable rather than merely discouraged. Third, the
honesty-immunity argument: the observation, made structural, that a benchmark which
makes no accuracy claim carries zero \Validated{} tags by design and is therefore
a stable kind of reference object---one that cannot be wrong about accuracy because
it asserts nothing about accuracy. The scoring taxonomy, the fault
detectability partition, and the monitors scored are all \Cited{}; we claim the
composition and the invariant, not the ingredients.

\paragraph{Limits, and what is deferred.}
The limits are deliberate and they set up the rest of the arc. The fault menu is
\Modelled{}: representative and locked by test on its honesty-critical structure,
but not exhaustive and not a stochastic model of any real network or clock. The
demonstration runs at $\IR=\num{e-3}$ on $64$-epoch scenarios; deep-tail integrity
rates such as \num{e-7} would require far more samples than any illustrative menu
carries, and resolving them is future work, not a claim made here. Real-capture
scenarios---driving the scorer with recorded rather than generated true-error
series---are likewise future work. And the benchmark scores monitors expressed as a
protection-level function; a monitor whose output is not reducible to a per-epoch
$\PL(\cdot)$ (a purely cryptographic freshness assertion with no timing bound, for
instance) can be framed by the menu but not fully scored by the coverage verdict, a
boundary we note rather than paper over. These are the seams the companion papers
of the arc are built to fill: the method paper \citep{baweja2026cti} supplies the
reference monitor demonstrated here, and the adversarial and common-mode extensions
supply the monitors whose boundaries the demonstration
(Section~\ref{sec:demo}) exposed.

\paragraph{Reading the numbers.}
Every magnitude in this paper---the $\SI{27}{\nano\second}$ handover floor, the
$\SI{200}{\nano\second}$ static delay, the \HMI{} counts, the availability
percentages---is a \Modelled{} number for a representative monitor and a
representative menu, regenerated by one command from the open engine. None is an
operational tolerance, a certified alert limit, or an endorsement. They illustrate
the yardstick; the yardstick, not the numbers, is what the \Proven{} invariant and
the honesty-immunity argument are about.
